\clearpage
\section{The parallel Fable experiment: preliminary review and comparison}
\label{app:parallel-fable}

\subsection{Scope and common setting}

Kinyon and Urban also ran Anthropic's Claude Fable~5.1 through Claude
Code on 13--15 September 2026~\cite{fable-kourovka}. The organizers
confirm that this used the same machine and experimental conditions:
the Kourovka Notebook, a 48-hour task, access to GAP and further software,
and requested local limits of generally 20 CPU cores and 100\,GB RAM.
This was a separate research run, later than the Codex run of
10--12 September. Fable is distinct from the Claude Opus[1m] agent that
reviewed this paper.

The comparison below was prepared by Codex on 17 September from the
29-page Fable paper and selected research reports and programs at
\href{https://github.com/JUrban/kour1cl/tree/290fc829d2be01cab8a1784d006bfc0250d7f11c}{repository commit \texttt{290fc829d2be}}.
It combines a preliminary reading of the arguments with three targeted
computational checks. It is smaller in scope than the external review
reported in Section~\ref{sec:external-review}; it does not establish
priority or certify the whole Fable portfolio. Configuration and large
search totals below are taken from that paper, rather than from a new
audit of its complete session.

\begin{center}\small
\begin{tabular}{@{}>{\raggedright\arraybackslash}p{0.21\textwidth}>{\raggedright\arraybackslash}p{0.35\textwidth}>{\raggedright\arraybackslash}p{0.35\textwidth}@{}}
\toprule
 & Codex experiment & Fable experiment\\\midrule
Research window (UTC) & 10 September 20:56 to 12 September 20:56 & 13 September 08:03 to 15 September 08:03\\[3pt]
Agent and harness & GPT-6 Astra; Codex; \texttt{xhigh} & Claude Fable~5.1; Claude Code~2.1.270; \texttt{xhigh} initially, then \texttt{high}\\[3pt]
LM delegation & One LM research agent & Six sub-agent calls for Notebook triage\\[3pt]
Reported answer entries & 46 deadline candidates; prior-work deductions listed separately & Five answer entries, including the main question of 2.78\\[3pt]
Further outcomes & Nine selected partial-result questions and bounded searches & Seven partial-result topics and a table of about 45 bounded verifications\\
\bottomrule
\end{tabular}
\end{center}

The bounded-verification table overlaps the Fable partial-result list,
so those counts should not be added. Nor are the two answer totals
normalized scores: 13.19 and 16.60 appear among Fable's five but outside
our 46, under prior-work status. The common local budget also does not
equate model inference expenditure, reasoning settings or delegation.
Fable records final consolidation at 09:44 on 15 September, about
101 minutes after its deadline, and reports that it added no new
mathematics then; that closeout included further verification and table
generation. A strict deadline comparison would need to separate those
activities in both archives.

\subsection{Reading the five answer entries}

\paragraph{21.32: finite derived-subgroup witnesses.}
Fable gives an explicit bound for a finite witness to \(G\cong H'\):
\[
 |H|\le |\Aut G|\,|Z(G)|^2
       \bigl(2\exp Z(G)\bigr)^{2\lceil\log_2|G|\rceil}.
\]
The argument bounds the index of a central subgroup and then replaces
that subgroup by its intersection with the derived subgroup. The
important step is the latter replacement: the universal coefficient
sequence and a Baer-sum argument preserve the isomorphism type of the
derived subgroup, not merely its order. The supplied proof addresses
this point explicitly and appears sound on this reading. With a
multiplication-table input, the bound yields decidability and a
nondeterministic quasi-polynomial upper bound. We have not undertaken
a literature search sufficient to determine priority for this result.

\paragraph{2.78: finiteness of two sets of group-theoretic numbers.}
Let \(f(G)\) count the subgroup orders that support a nonnormal subgroup.
Fable's construction gives
\[
 f\bigl(A_5\times M_{p^n}\times C_{p^r}\bigr)=7n+9r+9,
 \qquad p\ge7,\quad n\ge3,\quad r\ge0,
\]
where \(M_{p^n}\) is the modular \(p\)-group in its presentation.
The coprime-product formula and the subgroup argument for this family
give a convincing proof of the claimed finiteness. Our GAP spot checks
confirmed the family invariant in seven small parameter choices and
the subgroup-order count for \(A_5\).

The stated threshold \(k\ge86\) can be sharpened to \(k\ge78\)
without changing the construction. Taking \(n=3,\ldots,11\), the base
values \(7n+9\) are
\[
30,37,44,51,58,65,72,79,86.
\]
They represent every residue class modulo nine. In each class the
largest integer below its base value is at most \(86-9=77\), and
adding \(9r\) gives every subsequent value. Thus every \(k\ge78\)
occurs for a nonsoluble, nonsimple group. This small strengthening was
noticed during the present comparison. The further question in 2.78
about composite nonsoluble numbers, and the distribution question
3.57, remain open in the Fable report. Its additional classification
of the values \(0,\ldots,8\) uses deeper group theory and was not
independently verified in this review.

One statement in that classification needs the hypothesis that \(G\)
is nonsoluble: the assertion \(f(G)=7\) if and only if \(G\cong A_5\)
is false without it. The soluble dihedral group of order 512 has
nonnormal subgroups of exactly the seven orders
\(2,4,8,16,32,64,128\), as also confirmed by our check.
The Fable proof already assumes nonsolubility at this step; adding
it to the statement preserves the intended conclusion about
absolutely simple numbers.

\paragraph{18.46: a sharp square-root overgroup bound.}
The dihedral group \(D_8\) of order eight needs an overgroup of order
at least \(128=2|D_8|^2\) in which all its elements are squares.
We independently programmed a search for generators \(r,s\) with
\(|r|=4\), \(|s|=2\), \(r^s=r^{-1}\) and all eight elements of
\(\langle r,s\rangle\) in the square set. Unlike Fable's subgroup
enumeration, this uses element conjugacy classes and multiplication.
It excludes all 880 catalogue groups whose orders are
multiples of eight below 128, reproduces exactly the 14 successful
groups of order 128, and checks the witness \(D_8\wr C_2\).
This confirms the finite calculation, conditional on GAP and the
SmallGroups catalogue; the two implementations share that dependency.
The diagonal embedding in the wreath product also proves the upper
bound directly, since \(((g,1)\sigma)^2=(g,g)\), where \(\sigma\)
interchanges the two factors.

\paragraph{13.19: subdirect quotients.}
Fable constructs a group \(Q\) of order 512, with a normal subgroup
\(H\) of order 64 and three normal subgroups \(N_i\) of order 16,
such that \(\bigcap_iN_i=1\), \(HN_i=Q\) and \(Q/H\cong D_8\).
We rebuilt the multiplication from its bilinear formula over \(\F_2\)
in a separate Python program, checked the cocycle identity on all
\(64^3\) triples, and verified these subgroup and quotient properties.
Our example in Section~\ref{prior:13.19} has smaller \(Q\), of order
256, but uses four factors; Fable's uses three. These are different
size tradeoffs.

Both papers also give universal constructions for finite \(p\)-group
quotients. Fable's tagged-generator argument uses \(c+1\) factors for
a quotient of nilpotency class \(c\), whereas ours uses \(|P|\)
factors for a quotient \(P\). That factor bound is a useful additional
feature. The negative answer to the original regularity question
already follows from Kearnes--Mayr--Ru\v{s}kuc
\cite{kearnes-mayr-ruskuc}, as explained in
Section~\ref{prior:13.19}; this prior implication does not settle
the novelty of Fable's particular construction or factor bound.

\paragraph{16.60: the character-degree bound.}
The reduction to \(C_G(\alpha^2)\), followed by the twisted
Frobenius--Schur formula and monotonicity of the sum of character
degrees on subgroups, is a short, sound argument. An affirmative
answer was already stated in May 2026 by
Yerrapati--Dixit--Shukla~\cite[Theorem~3.1]{yerrapati-dixit-shukla}.
This is a rediscovery at the level of the problem's answer, with a
proof worth retaining; see also Section~\ref{prior:16.60}.

\subsection{Overlap, differences and useful corrections}

The overlapping questions show progress in both directions.
For 19.56, Fable excludes minimal simple groups and leaves a Frattini
lifting step; our proof in Section~\ref{cand:19.56} handles the
general finite-group question. For 20.21, Fable reports the stronger
lower bound \(|G|\ge6144\), compared with our deadline bound
\(|G|\ge3072\). We read its structural reductions but did not rerun
the large catalogue exclusions behind the stronger bound. For 20.100,
its finite search covers \(n=4\) in 85 groups of order at most 300
coprime to 30; Section~\ref{partial:20.100} gives the torsion-free
argument, the finite abelian reduction and certificates for
\(n=4,5,6,7\), with the previously stated public-artifact limitation
for \(n=7\). Fable also develops substantial partial work on 16.14
and 20.37 that is not part of our selected mathematical portfolio.

Two small corrections would improve the Fable exposition. Its
Appendix~B definition of \(L(G)\) for 20.37 should include
\(|A||B|=|G|\); without that condition, taking \(B=G\) makes the
definition trivial. In the proof for 21.111, the normalizer of a
cyclic subgroup generated by a unipotent element is \emph{contained
in} the indicated Borel subgroup, generally not equal to it. The
containment suffices for the subsequent involution argument. These
are local corrections, rather than objections to the stated results.

The two archives illustrate complementary research choices under
the same local constraints: Fable emphasizes finite computations and
constructions, while Codex's larger candidate list includes more
structural arguments and reuse of literature. The Fable results on
21.32, 2.78 and 18.46 deserve further review on their own merits.
A useful comparison should match exact problem parts, prior status
and verification scope, rather than interpret the headline totals
as a model ranking. Neither the shared organizers nor this later
cross-reading establishes a blind experimental comparison.

The pinned source manifest, new check programs and their outputs
accompany this paper in
\begin{center}\small
\artifact{paper/reviews/parallel-fable-2026-09-17/}.
\end{center}
The large Fable searches are available through its referenced repository.
